\documentclass[10pt,a4paper]{amsart}
\usepackage[margin=19mm]{geometry}
\usepackage{amsmath,amssymb,mathtools}
\usepackage[scr=rsfs,cal=euler]{mathalfa}
\usepackage{xfrac}
\usepackage[unicode,hidelinks,bookmarks=false]{hyperref}
\newcommand{\ie}{i.e.\ }
\newcommand{\cf}{cf.\ }
\newcommand{\resp}{respectively\ }
\newcommand{\Subsection}{\subsection}
\providecommand{\nobreakdash}{\nobreak\hbox{-}}
\setlength{\emergencystretch}{2em}
\multlinegap0pt
\usepackage{amssymb}
\usepackage{float}
\usepackage{enumerate}
\usepackage[normalem]{ulem}
\usepackage[shortlabels]{enumitem}
\newtheorem{theorem}{Theorem}
\title{Shock formation for 3D steady supersonic flows with general short pulse data}
\author{Bingbing Ding, Zhouping Xin, Huicheng Yin}
\date{Source: arXiv:2608.24599v1 (CC BY 4.0)}
\begin{document}
\maketitle

\noindent\emph{Source and licence.} Shock formation for 3D steady supersonic flows with general short pulse data. Version arXiv:2608.24599v1 (CC BY 4.0), distributed by arXiv under the Creative Commons Attribution 4.0 International licence, http://creativecommons.org/licenses/by/4.0/, as recorded in the arXiv licence metadata for this exact version and shown on the abstract page https://arxiv.org/abs/2608.24599v1. Full licence notice: \texttt{licenses/arxiv-2608.24599-CC-BY-4.0.txt}.\par\smallskip

\noindent\textit{Editorial scope.} Counted unit: the article's theorem Local (source0.tex lines 922--939). The article's complete original proof of it is delivered with it (source0.tex lines 941--1050, one \texttt{proof} environment of 110 lines). No mathematics is written, repaired or re-derived: the statement and the proof are the article's own lines, in the article's own notation, and no retained line is substituted or re-flowed.  Retained uncounted as the source-local context the proof needs: lines 638--644.  Nothing was omitted from any retained span, so the package carries no omission record.  No retained line contains a citation, so the wrapper carries no bibliography and no bibliography entry.  No retained line contains a figure environment, an image-inclusion command or drawing code, so no image is delivered and the package declares no asset dependency.  Editorial formatting only, in addition: an AMS \texttt{amsart} preamble that restates the article's own declarations and macros used by the retained lines, namely the environments \texttt{theorem} on separate counters, together with the standard packages those lines need.  The retained environments print in this document as Theorem 1 in order of appearance, whereas the article prints the same items with the numbering of its own full document.  The complete native source of the article as distributed in the arXiv e-print for this exact version is bundled unchanged as \texttt{sources/208493/source0.tex}.
\begin{equation}\label{mainequation}
\begin{cases}
\displaystyle\sum_{\alpha,\beta=0}^2 g^{\alpha\beta} \partial_{\alpha\beta}^2 \phi = 0,\\[6pt]
\phi|_{t=0}=\delta^{\nu}\Phi_0(\frac{\mathcal{A}r - 1}{\delta}, \omega), \quad
\partial_{0}\phi|_{t=0}=\delta^{\nu - 1}\Phi_1(\frac{\mathcal{A}r - 1}{\delta}, \omega),
\end{cases}
\end{equation}

\begin{theorem}\label{Local}
For small $\delta > 0$, \eqref{mainequation} admits a local solution $\phi \in C^{\infty}([0, t_0] \times \mathbb{R}^2)$.
Meanwhile, for any $m \in \mathbb{N}_0$, $q \in \mathbb{N}_0^3$ and $k \in \mathbb{N}_0$, the following estimates hold
\begin{align}
&|L^m \partial^q \Omega^k \phi(t_0, y)|
\leq \delta^{\nu-m(2-\nu)-|q|} \quad \text{for } \frac{1 - 2\delta}{\mathcal{A}} \leq r \leq \frac{1 + 2\delta}{\mathcal{A}},\label{initial1}\\
&|\underline L^m \partial^q \Omega^k \phi(t_0, y)|
\leq \delta^{\nu-m(2-\nu)-|q|} \quad \text{for } \frac{1 - 3\delta}{\mathcal{A}} \leq r \leq \frac{1 + \delta}{\mathcal{A}}.\label{initial2}
\end{align}
In addition, for any $r_0\in (\frac{1-3\delta}{\mathcal A}, \frac{1+2\delta}{\mathcal A})$,
\begin{equation}\label{y0}
\delta^{2-\nu}\underline L\partial_{0}\phi(t_0, r_0\omega)=\mathcal{A} \big(\mathcal A\partial_s^2\Phi_0\big(\tfrac{\mathcal A(r_0-t_0)-1}{\delta}, \omega\big)-\partial_s\Phi_1\big(\tfrac{\mathcal A(r_0-t_0)-1}{\delta}, \omega\big)\big)+ O(\delta^{\nu-1})
\end{equation}
and
\begin{equation}\label{y1}
\delta^{2-\nu}L\partial_{0}\phi(t_0, r_0\omega)=\mathcal{A} \big(\mathcal A\partial_s^2\Phi_0\big(\tfrac{\mathcal A(r_0+t_0)-1}{\delta}, \omega\big)+\partial_s\Phi_1\big(\tfrac{\mathcal A(r_0+t_0)-1}{\delta}, \omega\big)\big)+ O(\delta^{\nu-1}).
\end{equation}
\end{theorem}

\begin{proof}
At first, we give the following bootstrap assumption:
\begin{equation}\label{ba}
\begin{split}
&\text{\it when $0\leq t\leq t_0$ and integer $N_0\ge 6$,
$|\partial^\kappa\Omega^k\phi|\leq\delta^{\nu_0-|\kappa|}$ holds for $|\kappa|+k\leq N_0$},\\
&\text{\it where $\nu_0\in(1, \nu)$ is any fixed constant.}
\end{split}
\end{equation}
For \eqref{mainequation} and $n\in\mathbb N_0$, define the energy
$$
M_n(t)=\sum_{|\kappa|+k\leq n}\delta^{2|\kappa|}\|\partial\partial^\kappa\Omega^k\phi(t,\cdot)\|_{L^2(\mathbb R^2)}^2.
$$
Set $w=\delta^{|\kappa|}\partial^\kappa\Omega^k\phi$ with $|\kappa|+k\leq 2N_0-2$.
It follows from \eqref{mainequation} and integration by parts that
\begin{equation}\label{Y-4}
\begin{split}
&\int_0^t\int_{\Sigma_\tau}\partial_t w\,g^{\alpha\beta}\partial_{\alpha\beta}^2 w\,dyd\tau\\
=&\int_{\Sigma_{t}}\big(g^{0\alpha}\partial_\alpha w\,\partial_t w
-\frac12g^{\alpha\beta}\partial_\alpha w\,\partial_\beta w\big)dy
-\int_{\Sigma_{0}}\big(g^{0\alpha}\partial_\alpha w\,\partial_t w
-\frac12g^{\alpha\beta}\partial_\alpha w\,\partial_\beta w\big)dy\\
&+\int_0^t\int_{\Sigma_\tau}\big(-(\partial_\alpha g^{\alpha\beta})\partial_\beta w\,\partial_t w
+\frac12(\partial_t g^{\alpha\beta})\partial_\alpha w\,\partial_\beta w\big)dyd\tau,
\end{split}
\end{equation}
where
\begin{equation}\label{w}
\begin{split}
|g^{\alpha\beta}&\partial_{\alpha\beta}^2w|
\lesssim\delta^{|\kappa|}\sum_{\substack{
|\kappa_1|+|\kappa_2|\leq|\kappa|,\ k_1+k_2\leq k,\\
|\kappa_2|+k_2+2<|\kappa|+k}}
|\partial^{\kappa_1}\Omega^{k_1}g^{\alpha\beta}|
\cdot|\partial^{\kappa_2}\partial^2\Omega^{k_2}\phi|.
\end{split}
\end{equation}
Combining \eqref{ba}--\eqref{w} with Gronwall's inequality yields that for $0\leq t\leq t_0$,
$$
M_{2N_0-2}(t)\lesssim M_{2N_0-2}(0)\lesssim\delta^{2\nu-1}.
$$
Applying the Sobolev embedding on the circle $\mathbb S_r$ (centered at the origin with radius $r$), we arrive at
$$
|w(t,y)|\lesssim \frac{1}{\sqrt{r}}\|\Omega^{\le 1}w\|_{L^2(\mathbb S_r)}.
$$
Due to $r\sim 1$ for $t\in[0,t_0]$ and $(t,y)\in\mathrm{supp}\,w$, then one has that for $|\kappa|+k\leq N_0$,
\begin{equation}\label{local}
|\partial^\kappa\Omega^k\phi(t,y)|
\lesssim \|\Omega^{\le 1}\partial^\kappa\Omega^k\phi\|_{L^2(\mathbb S_r)}
\lesssim\delta^{1/2}\|\partial\Omega^{\le 1}\partial^\kappa\Omega^k\phi\|_{L^2(\Sigma_{t})}
\lesssim\delta^{\nu-|\kappa|}.
\end{equation}
Therefore, for small $\delta>0$ and $\nu_0<\nu$, the bootstrap assumption \eqref{ba} can be closed by the continuity argument.


Based on \eqref{local}, we now improve the $L^\infty$ estimate of $\phi$ on the domain $D_1=\{(t,y): 0\le t \le t_0,$
$\frac{1}{\mathcal A} - t\le r \le t+\frac{ 1}{\mathcal A}\}$.
The equation in \eqref{mainequation} can be rewritten as
\begin{equation}\label{me}
L\underline{L}\phi
= \frac{1}{2r}L\phi
- \frac{1}{2r}\underline{L}\phi
+ \frac{1}{r^2}\Omega^2\phi
+ 2g^{0i}\partial_{0i}^2\phi
+ (g^{ij}-m^{ij})\partial_{ij}^2\phi,
\end{equation}
where $m^{ii} = 1$ for $i=1, 2$ and $m^{ij} = 0$ for $i\neq j$.
It follows from \eqref{local} and \eqref{me} that
\begin{equation}\label{underlineLL}
|\underline{L}L\partial^q\Omega^k\phi(t,y)|
\le \delta^{2\nu-3-|q|},
\quad 0 \le t \le t_0.
\end{equation}
Note that $\phi$ vanishes on the surface $\{(t,y) : r- t = \frac{1}{\mathcal A}\}$, then integrating \eqref{underlineLL} along
the integral curves of $\underline{L}$ yields
\begin{equation}\label{YHCC4}
|L\partial^q\Omega^k\phi(t,y)|
\le \delta^{2\nu-2-|q|},
\quad (t,y) \in D_1.
\end{equation}
By applying $L$ repeatedly to both sides of \eqref{me} and using an induction argument,
we conclude that for $m\in\Bbb N$,
\[
|L^m\partial^q\Omega^k\phi(t,y)|
\le \delta^{\nu-m(2-\nu)-|q|},
\quad (t,y) \in D_1,
\]
which yields \eqref{initial1}.


Similarly, for any $(t,y) \in D_2=\{(t,y): 0\le t \le t_0,
\frac{1-\delta}{\mathcal A} - t \le r \le t+\frac{ 1-\delta}{\mathcal A}\}$,
\[
|\underline L^m\partial^q\Omega^k\phi(t,y)|
\le \delta^{\nu-m(2-\nu)-|q|},
\]
and hence \eqref{initial2} is proved.

For any fixed $\omega\in\Bbb S$, since both $(0,(r_0-t_0)\omega)$ and $(t_0, r_0\omega)$ lie on the line $\{(t,r\omega): r-t=r_0-t_0\}$,
then it follows from \eqref{me} that
\begin{equation}\label{YHCCC-3}
\begin{split}
&\underline{L}\partial_{0}\phi(t_0,r_0\omega)
=\underline{L}\partial_{0}\phi(0,(r_0-t_0)\omega)+\int_{(0,(r_0-t_0)\omega)}^{(t_0,r_0\omega)}L\underline{L}\partial_{0}\phi\\
=&\delta^{\nu-2}\mathcal{A}\big(\mathcal{A}\partial_s^2\Phi_0\big(\frac{\mathcal{A}(r_0-t_0)-1}{\delta}, \omega\big)-\partial_s\Phi_1\big(\frac{\mathcal{A}(r_0-t_0)-1}{\delta}, \omega\big)\big)+O(\delta^{2\nu-3}),
\end{split}
\end{equation}
which yields \eqref{y0}. Analogously, \eqref{y1} can be obtained
by integrating \eqref{me} along the integral curves of $\underline L$.
\end{proof}

\end{document}
